\begin{abstractpage}

\begin{abstract}{romanian}
În contextul actual, înțelegerea informațiilor nutriționale de pe etichetele alimentelor reprezintă o provocare constantă pentru consumatori, din cauza terminologiei complexe și a listelor lungi de ingrediente cu denumiri științifice. Această lucrare prezintă dezvoltarea unei aplicații pentru recunoașterea și interpretarea automată a informațiilor nutriționale folosind tehnologii avansate de recunoaștere optică a caracterelor și procesarea limbajului natural.

Aplicația LabelLogic transformă listele de ingrediente cu denumiri științifice într-un format ușor de înțeles, oferind explicații despre rolul fiecărui component și calculând un scor de procesare pentru evaluarea gradului de procesare a alimentului. Soluția oferă funcționalități de filtrare și comparare a alimentelor după diverse criterii nutriționale, facilitând utilizatorilor luarea unor decizii informate despre alimentație.

Rezultatele demonstrează fezabilitatea utilizării tehnologiilor contemporane pentru automatizarea analizei nutriționale, oferind o soluție practică pentru promovarea unei alimentații mai sănătoase.
\end{abstract}

\begin{abstract}{english}
In the current context, understanding nutritional information on food labels is a constant challenge for consumers due to complex terminology and long lists of ingredients with scientific names. This paper presents the development of an application for automatic recognition and interpretation of nutritional information using advanced optical character recognition and natural language processing technologies.

The application transforms ingredient lists with scientific names into an easy-to-understand format, providing explanations of the role of each component and calculating a processing score for evaluating the degree of food processing. The solution provides functionality to filter and compare foods according to various nutritional criteria, making it easier for users to make informed dietary decisions.

The results demonstrate the feasibility of using contemporary technologies to automate nutritional analysis, providing a practical solution for promoting healthier diets.
\end{abstract}

\end{abstractpage}